problem:  
Let \((M^n,g)\) be a compact Riemannian manifold with \(\operatorname{Ric}_g \ge (n-1)k\, g\) for some \(k>0\).  
Let \(I_M(v)\) denote the *isoperimetric profile function*, i.e.  
\[
I_M(v)=\inf\{\operatorname{Area}(\partial\Omega):\operatorname{Vol}(\Omega)=v\}, \quad v\in(0,\operatorname{Vol}(M,g)).
\]  
The Lévy–Gromov inequality asserts that  
\[
I_M(v) \ge I_{S_k^n}(v),
\]  
where \(S_k^n\) is the sphere of constant sectional curvature \(k\), and equality for all \(v\) characterizes the round sphere itself.

**Conjecture (Partial Equality Rigidity):**  
Assume that there exist two distinct nontrivial volumes \(v_1,v_2\in(0,\operatorname{Vol}(M,g))\), with \(v_1\neq v_2\), such that  
\[
I_M(v_i)= I_{S_k^n}(v_i),\quad i=1,2,
\]  
and the corresponding isoperimetric regions realizing these values have smooth boundaries (for example, if \(n\le 7\)).  
Does it follow that \((M,g)\) is isometric to the constant–curvature sphere \(S_k^n\)?

Why is it a "good" problem:  
This conjecture explores a subtle and challenging *rigidity from partial equality* phenomenon at the interface of metric geometry, geometric analysis, and Ricci curvature comparison theory. The statement is minimal and elegant—it only modifies the fully sharp case of Lévy–Gromov by asking whether equality at two nontrivial volumes already enforces complete geometric rigidity—but addressing it likely requires new insight into the global structure of isoperimetric profiles.

The depth of the problem lies in the following:

1. **Concavity under Ricci bounds:** Under \(\operatorname{Ric}\ge (n-1)k\), the normalized isoperimetric profile satisfies a concavity-type differential inequality. If equality at two points forces equality of the differential inequality, one might deduce that it holds identically and thus recover the full spherical model. Analyzing this rigorously involves delicate control of second variations and stability of minimizers—an open and technically profound challenge.

2. **Geometric measure theory meeting curvature comparison:** Understanding when local equality of isoperimetric data propagates to global equality touches on the regularity of isoperimetric hypersurfaces, second fundamental form estimates, and measure-theoretic curvature bounds—bridging tools from analysis, PDE, and Riemannian geometry.

3. **Cross-disciplinary implications:** Any resolution would deepen our understanding of rigidity in comparison geometry, influence developments in synthetic curvature bounds (in metric measure spaces à la Lott–Sturm–Villani), and might suggest new mechanisms by which partial equality information determines full geometric structure.

4. **Extreme simplicity and profundity:** The statement involves only two equality points of a well-known geometric function but has potentially sweeping consequences for our understanding of curvature–geometry rigidity principles.

Hence, this question embodies the essential features of a *good* problem: it is simple to state, deeply rooted in current research frontiers, connects multiple fields, and, if resolved, would significantly advance geometric analysis and our comprehension of curvature-induced rigidity.